\chapter{Primitive content and conserved arithmetic sectors}\label{ch:content}
\section{The coordinate gcd is a real invariant}
For $v\in\Vn$, let $g(v)$ be the positive gcd of its five coordinates. Since the coordinates occupy all residues modulo $5$, one has $5\nmid g(v)$. Call $g(v)$ the content of the packet. A primitive packet has content $1$.

The content is compatible with the shell in a precise way. If $g=g(v)$, define
\begin{equation}\label{eq:contentscale}
u_j=\frac{v_{gj\bmod5}}g,\qquad j=0,\ldots,4.
\end{equation}
The coordinate permutation restores the standard residues after division. The vector $u$ has sum zero and gcd $1$. Its norm squared is $10n/g^2$. For every integral standard-residue vector of sum zero, the norm is divisible by $5$, because its squared residues sum to $30$, and by $2$, because squares and coordinates agree modulo $2$. Thus $10n/g^2$ is $10$ times an integer, and
\[
g^2\mid n.
\]
The map is a bijection from exact-content-$g$ packets on shell $n$ to primitive packets on shell $n/g^2$.

The permutation $j\mapsto gj\bmod5$ has sign $\chifive(g)$. For the four units modulo $5$, multiplication by $1$ is the identity, multiplication by $2$ or $3$ is a four-cycle on the nonzero residues, and multiplication by $4$ is two transpositions. Scaling the degree-$10$ Vandermonde therefore gives
\begin{equation}\label{eq:wscale}
W(v)=\chifive(g)g^{10}W(u).
\end{equation}
The character is not an arbitrary normalization. It records the sign of the residue-restoring permutation.

\section{The primitive residual and its inverse}
Define
\begin{equation}\label{eq:residualdef}
R(n)=\sum_{\substack{v\in\Vn\\g(v)=1}}W(v).
\end{equation}
Partitioning the packets by content and using \eqref{eq:wscale} yields
\begin{equation}\label{eq:contentdecomposition}
\tau(n)=\sum_{g^2\mid n}\chifive(g)g^{10}R(n/g^2).
\end{equation}
Terms with $5\mid g$ have character zero and correspond to no actual content sector. They may be included in the formula without adding contributions.

The square-supported arithmetic kernel in \eqref{eq:contentdecomposition} has inverse
\[
d^2\longmapsto\mu(d)\chifive(d)d^{10}.
\]
Indeed, at a square index $m^2$, the convolution of the two kernels is
\[
\chifive(m)m^{10}\sum_{d\mid m}\mu(d),
\]
which is $1$ at $m=1$ and $0$ otherwise. At nonsquare indices both kernels have no support. Consequently
\begin{equation}\label{eq:mobius}
\boxed{R(n)=\sum_{d^2\mid n}\mu(d)\chifive(d)d^{10}\tau(n/d^2).}
\end{equation}
Both formulas are exact for every positive integer. They are not fitted from the first few shells.

At $n=1$, $R(1)=1$. At a prime, only $d=1$ occurs in \eqref{eq:mobius}, giving
\begin{equation}\label{eq:rprime}
R(p)=\tau(p).
\end{equation}
The primitive extraction has no radial contribution to remove at a prime. This is the simplest reason it cannot, by itself, settle the prime nonvanishing obligation.

\section{Homogeneous universal nulls preserve each content}
\begin{theorem}\label{thm:homcontent}
Let $Q$ be a symmetric homogeneous polynomial of degree $d$ that is universally null. Then its weighted sum is zero separately in every exact-content sector on every shell.
\end{theorem}
\begin{proof}
Let
\[
R_Q(n)=\sum_{\substack{v\in\Vn\\g(v)=1}}W(v)Q(v).
\]
Symmetry removes the residue permutation from $Q$, and homogeneity gives $Q(v)=g^dQ(u)$. The scaling bijection yields
\begin{equation}\label{eq:thetacontent}
\sum_{v\in\Vn}W(v)Q(v)
=\sum_{g^2\mid n}\chifive(g)g^{10+d}R_Q(n/g^2).
\end{equation}
The left side is zero. At $n=1$, only $g=1$ contributes, so $R_Q(1)=0$. Inductively, the $g=1$ term is $R_Q(n)$ and every nonzero term with $g>1$ has smaller argument. Thus $R_Q(n)=0$ for every $n$. The contribution of exact content $g$ is its scaling factor times $R_Q(n/g^2)$, so it is zero as well.
\end{proof}

Finite linear combinations of homogeneous universal nulls inherit this conservation law even when their degrees differ. Chapter~\ref{ch:pcall} will prove the stronger result for every finite symmetric polynomial that is universally null, without initially requiring a decomposition into known homogeneous nulls.

Parity sectors failed to conserve general null transfers. Content sectors do conserve them. This is the distinction between a geometric partition chosen for convenience and an arithmetic decomposition forced by a scaling identity.

\section{A composite obstruction with opposite conserved masses}
At $n=121$, the contents are $1$ and $11$. The exact masses are
\[
R(121)=-25439104668,\qquad11^{10}R(1)=25937424601.
\]
Their sum is
\[
\tau(121)=498319933.
\]
One can verify these values independently from $\tau(11)=534612$:
\[
\tau(121)=\tau(11)^2-11^{11},\qquad
R(121)=\tau(121)-11^{10}.
\]
The two content masses have opposite nonzero signs. Every retained homogeneous universal-null correction preserves them separately. Hence no correction in that entire permitted span can make all real corrected packet weights have one common sign at $121$. Increasing the degree cannot remove this obstruction.

The obstruction is composite-specific. A prime shell has content $1$ only. There the conservation law adds no invariant beyond the total. It must not be promoted into an impossibility theorem for prime-specific sign constructions.

\figplaceholder{content121}{The two exact-content masses at $121$ have opposite signs. Their sum survives, but every permitted finite symmetric polynomial universal null conserves the two masses separately. A common-sign correction at this composite shell is therefore impossible in that class.}

\section{A rank bound from content}
Suppose a shell has $N$ $J$-orbits occupying $c$ nonempty content classes. A transfer vector produced by a homogeneous universal null has zero sum in each class. That subspace has dimension $N-c$, so every such transfer matrix has rank at most $N-c$.

The constructed family through degree $28$ attains this bound at selected composite shells:
\begin{center}
\begin{tabular}{rrrr}
\toprule
$n$&$J$-orbits&Content classes&Rank\\
\midrule
9&4&$1,3$&2\\
49&22&$1,7$&20\\
81&31&$1,3,9$&28\\
\bottomrule
\end{tabular}
\end{center}
These are exact finite ranks. They illustrate the conserved split and the effectiveness of the constructed transfers. They do not prove saturation of the content-preserving space at every shell.

\chapter{Refinement and exact transfer spaces}\label{ch:transport}
\section{Mass vectors and aggregation}
Fix a shell, a finite rational family $Q_1,\ldots,Q_h$ of universally null corrections, and the occupied sectors at depth $k$. Let their number be $s_k$. Define
\[
t_{k,r}=\sum_{v\in S_{k,r}}W(v),\qquad
(M_k)_{rj}=\sum_{v\in S_{k,r}}W(v)Q_j(v).
\]
The vector $t_k$ contains sector masses. The columns of $M_k$ are the effects of the permitted corrections. Coefficients $a\in\Q^h$ change the mass vector to $t_k+M_ka$.

Every depth-$k+1$ sector belongs to exactly one parent at depth $k$. Let $A$ be the corresponding $0$--$1$ aggregation matrix. Summing child contributions gives the exact identities
\begin{equation}\label{eq:aggregate}
t_k=At_{k+1},\qquad M_k=AM_{k+1}.
\end{equation}
Every occupied parent has a child. Each correction column has global total zero.

Put
\[
H_k=\ker(\mathbf1^T:\Q^{s_k}\to\Q),\qquad
T_k=\im M_k,\qquad E_k=H_k/T_k.
\]
The quotient $E_k$ is the vector space of total-preserving changes that are unavailable from the chosen family. It is not a count of failed sign certificates. A sign certificate may exist even when $E_k$ is nonzero.

\section{The exact refinement sequence}
\begin{theorem}\label{thm:refinement}
There is a short exact sequence
\begin{equation}\label{eq:refinement}
0\longrightarrow
\frac{\ker A}{M_{k+1}(\ker M_k)}
\longrightarrow E_{k+1}\longrightarrow E_k\longrightarrow0.
\end{equation}
Consequently, if $r_k=\rank M_k$,
\begin{equation}\label{eq:defectgrowth}
\dim E_{k+1}-\dim E_k
=(s_{k+1}-s_k)-(r_{k+1}-r_k).
\end{equation}
\end{theorem}
\begin{proof}
The aggregation map sends $H_{k+1}$ onto $H_k$. Given a zero-total parent vector, place each parent entry in one selected child and zero in its other children. The resulting child vector has zero total and maps to the required vector.

Equation \eqref{eq:aggregate} gives $A(T_{k+1})=T_k$, so aggregation induces a surjective map between the quotients. Its kernel is the quotient of $\ker A$ by $T_{k+1}\cap\ker A$. A transfer $M_{k+1}a$ lies in $\ker A$ exactly when
\[
AM_{k+1}a=M_ka=0.
\]
Thus $T_{k+1}\cap\ker A=M_{k+1}(\ker M_k)$, which proves \eqref{eq:refinement}. Finally $\dim E_k=s_k-1-r_k$, giving \eqref{eq:defectgrowth}.
\end{proof}

This is ordinary exact linear algebra applied to the packet data. Refinement introduces $s_{k+1}-s_k$ new detail directions. Effective new column rank resolves $r_{k+1}-r_k$ of them. The formula says exactly how much of the new geometric detail remains unavailable.

\section{Localization by cancellation, rather than indicators}
Let the columns of a rational matrix $N$ form a basis of $\ker M_k$. Then
\begin{equation}\label{eq:detailmatrix}
B_k=M_{k+1}N,\qquad AB_k=0.
\end{equation}
The columns of $B_k$ transfer mass only within the parents: their effects on every parent cancel. They are still combinations of universal null polynomials. This is the valid repair of the failed sector-indicator proposal.

There is a cost. The finer effects must be computed and solved for, and their span may be smaller than all child changes preserving parent masses. Equation \eqref{eq:refinement} identifies that remaining detail defect. One does not obtain local freedom simply by declaring that a global null polynomial is available in each sector separately.

\section{The exact calculation at 31}
The shell at $31$ has six coarse occupied sectors and sixteen terminal $J$-orbit sectors. The constructed ranks are
\begin{center}
\begin{tabular}{rrrrr}
\toprule
Degree cutoff&Coarse rank&Orbit rank&Orbit defect&Detail rank\\
\midrule
14&5&5&10&0\\
20&5&14&1&9\\
22&5&15&0&10\\
\bottomrule
\end{tabular}
\end{center}
The zero-total coarse space has dimension $5$, and the zero-total orbit space has dimension $15$. At degree $22$, the available columns fill both. The detail space has dimension $15-5=10$.

The stored two-stage certificate first eliminates five coarse masses, using only five elimination equations. It then uses coefficient vectors in $\ker M_1$ to eliminate all but one fine mass while preserving every coarse total. The survivor is
\[
-52843168.
\]
The product engine independently identifies this with $\tau(31)$. The solve does not require a preassigned tau value in its elimination equations. The final survivor remains an exact finite calculation, not a formula that independently forces a nonzero survivor at every prime.

The degree-$20$ family has an orbit-transfer defect of dimension $1$, but it already supplied a strictly negative sign certificate at $31$. Filling every zero-total direction is stronger than needed for that finite sign task. Conversely, filling them at degree $22$ does not by itself prove a nonzero total. These two facts prevent rank from being misinterpreted as a nonvanishing test.

\figplaceholder{ranks31}{The exact ranks at $31$. The coarse rank is already full at degree $14$, while the terminal orbit rank reaches $15$ at degree $22$. A degree-$20$ sign certificate exists despite its one-dimensional transfer defect. Rank and nonvanishing are different statements.}

\section{Other finite saturation profiles}
Full zero-total orbit rank is obtained in the retained family at
\begin{center}
\begin{tabular}{rrr}
\toprule
Prime&$J$-orbits&Degree cutoff\\
\midrule
31&16&22\\
37&19&22\\
41&22&24\\
59&30&26\\
71&36&28\\
\bottomrule
\end{tabular}
\end{center}
These cutoffs are not asserted optimal among all polynomials. The first five corrections through degree $14$ also had full coarse zero-total rank at all $94$ odd primes below $500$. Among the $250$ odd indices below $500$, the only coarse failures were $9$ and $81$. The content split explains genuine conserved directions in those composite examples.

The fresh transport replay verifies $144$ refinement ranks and the complete rational certificate at $31$. Its results authenticate these finite matrices. A uniform theorem about every prime, every degree requirement or every content-preserving transfer is not obtained by these rank profiles.

\section{What further completeness would mean}
A plausible geometric question is whether all-degree universal nulls span every content-preserving transfer at each shell. Chapter~\ref{ch:interpolation} explains why ordinary polynomial interpolation does not settle it. Even if such completeness were proved, the conserved total at a prime would remain $\tau(p)$.

The useful implication of completeness would be that every compatible redistribution is expressible by legitimate corrections. The missing arithmetic implication would still be that the compatible total is nonzero. Treating these as separate statements prevents a successful linear-algebra calculation from silently doing the work of the unsolved conjecture.
